\subsection{幺半群的融合和}

设 \((M_{i})_{i\in I}\) 表示一族幺半群，\(e_i\) 表示 \(\mathbf{M}_i\) 的单位元。给定一个幺半群 \(\mathbf{A}\) 以及一族同态 \(h_i\colon \mathbf{A}\to \mathbf{M}_i\) （\(i\in I\) ）。

集合 S 即族 \((\mathbf{M}_{i})_{i\in\mathbb{I}}\) 的和，其元素为有序对 \((i,x)\) ，其中 \(i\in I\) 与 \(x\in M_{i}\) 。对每个三元组 \(\alpha=(i,x,x^{\prime})\) ，其中 \(i\in I\) ，\(x,x^{\prime}\) 在 \(M_{i}\) 中，记 \(u_{\alpha}=(i,xx^{\prime})\) 与 \(v_{\alpha}=(i,x).(i,x^{\prime})\) ；对每个三元组 \(\lambda=(i,j,a)\) （属于 \(I\times I\times A\) ），记 \(p_{\lambda}=(i,h_{i}(a))\) 与 \(q_{\lambda}=(j,h_{j}(a))\) ；对所有 \(i\in I\) ，记 \(\varepsilon_{i}=(i,e_{i})\) 。由 S 和关系子 \((u_{\alpha},v_{\alpha})\) ， \((p_{\lambda},q_{\lambda})\) 与 \((\varepsilon_{i},e)\) 所定义的幺半群 M 称为族 \((\mathbf{M}_{i})_{i\in\mathbb{I}}\) 关于 A 的融合和。设 \(\phi\) 表示 \(\mathrm{Mo}(\mathrm{S})\) 到 M 上的典范同态，并记 \(\phi_{i}(x)=\phi(i,x)\) ，其中 \((i,x)\in\mathrm{S}\) 。称 \(\phi_{i}\) 为 \(M_{i}\) 到 M 中的典范映射。对所有 \(a\in A\) ，元素 \(\phi(i,h_{i}(a))\) 不依赖于 i，记作 \(h(a).†\)

由生成元和关系子定义的幺半群的泛性质（第 2 号）蕴含以下结果：

命题 4. (a) 对所有 \(i \in \mathbf{I}\) ，映射 \(\phi_i\) 是 \(\mathbf{M}_i\) 到 \(\mathbf{M}\) 中的同态，且 \(\phi_i \circ h_i = h\) 对所有 \(i \in \mathbf{I}\) 成立。进一步， \(\mathbf{M}\) 由 \(\bigcup_{i \in \mathbf{I}} \phi_i(\mathbf{M}_i)\) 生成。

\begin{enumerate}
\def\labelenumi{(\alph{enumi})}
\setcounter{enumi}{1}
\tightlist
\item
  设 \(\mathbf{M}'\) 是一个幺半群，且 \(f': \mathbf{M}_i \to \mathbf{M}'\) （\(i \in \mathbf{I}\) ）为一族同态，使得 \(f_i \circ h_i\) 与 \(i \in \mathbf{I}\) 无关。存在且仅存在一个同态 \(f: \mathbf{M} \to \mathbf{M}'\) ，使得 \(f_i = f \circ \phi_i\) 对所有 \(i \in \mathbf{I}\) 成立。
\end{enumerate}

在下文中，我们将作如下假设：

\begin{enumerate}
\def\labelenumi{(\Alph{enumi})}
\tightlist
\item
  对所有 \(i \in \mathbf{I}\) ，存在一个子集 \(\mathbf{P}_i\) ，它是 \(\mathbf{M}_i\) 的包含 \(e_i\) 的子集，使得映射 \((a, p) \mapsto h_i(a).p\) 从 \(\mathbf{A} \times \mathbf{P}_i\) 到 \(\mathbf{M}_i\) 是双射。
\end{enumerate}

这意味着诸同态 \(h_i\) 都是单射。设 \(x \in \mathbf{M}\) ；每个有限序列 \(\sigma = (a: i_1, \ldots, i_n; p_1, \ldots, p_n)\) ，其中 \(a \in \mathbf{A}\) , \(i_\alpha \in \mathbf{I}\) 且 \(p_\alpha \in \mathbf{P}_{i_\alpha}\) （\(1 \leqslant \alpha \leqslant n\) ），若满足

\[
x = h (a) \cdot \prod_ {\alpha = 1} ^ {n} \phi_ {i _ {\alpha}} (p _ {\alpha})\tag{3}
\]

便称为 \(x\) 的一个分解。整数 \(n \geqslant 0\) 称为分解 \(\sigma\) 的长度，记作 \(l(\sigma)\) ；序列 \((e)\) 是 \(M\) 的单位元的一个长度为 0 的分解。分解 \(\sigma\) 称为既约的，如果 \(i_{\alpha} \neq i_{\alpha + 1}\) 对 \(1 \leqslant \alpha < n\) 成立且 \(p_{\alpha} \neq e_{i_{\alpha}}\) 对 \(1 \leqslant \alpha \leqslant n\) 成立。

命题 5. 在假设 (A) 下，每个元素 \(x\) （属于 \(\mathbf{M}\) ）都有唯一的既约分解 \(\sigma\) 。每个分解 \(\sigma' \neq \sigma\) ，只要它分解的是 \(x\) ，就满足 \(l(\sigma') > l(\sigma)\) 。

\BourbakiRunInHeading{(A) 既约分解的唯一性：}

设 \(\Sigma\) 表示由满足如下条件的序列 \(\sigma = (a; i_1, \ldots, i_n; p_1, \ldots, p_n)\) 组成的集合： \(n \geqslant 0\) , \(a \in \mathbf{A}\) , \(i_\alpha \in \mathbf{I}\) 且 \(p_\alpha \in \mathrm{P}_{i_\alpha} - \{e_{i_\alpha}\}\) （\(1 \leqslant \alpha \leqslant n\) ），并且 \(i_\alpha \neq i_{\alpha+1}\) （\(1 \leqslant \alpha < n\) ）。设 \(\Phi\) 表示由下式定义的 \(\Sigma\) 到 \(\mathbf{M}\) 中的映射：

\[
\Phi (a; i _ {1}, \dots , i _ {n}; p _ {1}, \dots , p _ {n}) = h (a). \prod_ {\alpha = 1} ^ {n} \phi_ {i _ {\alpha}} (p _ {\alpha}).\tag{4}
\]

\(x \in \mathbf{M}\) 的一个既约分解是某个元素 \(\sigma\) ，它属于 \(\Sigma\) 且满足 \(\Phi(\sigma) = x\) 。对所有 \(i \in \mathbf{I}\) ，设 \(\Sigma_i\) 为 \(\Sigma\) 的由序列 \((e; i_1, \ldots, i_n; p_1, \ldots, p_n)\) （其中 \(i \neq i_1\) ，当 \(n > 0\) ）组成的子集。设

\[
\sigma = (e; i _ {1}, \dots , i _ {n}; p _ {1}, \dots , p _ {n})
\]

为 \(\Sigma_{i}\) 中的序列，并设 \(\xi\) 在 \(\mathbf{M}_i\) 中；令 \(\xi = h_i(a).p\) ，其中 \(a\in \mathbf{A}\) , \(p\in \mathbf{P}_i\) ，且

\[
\Psi_ {i} (\xi , \sigma) = \left\{ \begin{array}{l l} (a; i _ {1}, \dots , i _ {n}; p _ {1}, \dots , p _ {n}) & \text { if } p = e _ {i} \\ (a; i, i _ {1}, \dots , i _ {n}; p, p _ {1}, \dots , p _ {n}) & \text { if } p \neq e _ {i}. \end{array} \right.\tag{5}
\]

显然 \(\Psi_{i}\) 是 \(\mathbf{M}_i\times \Sigma_i\) 到 \(\Sigma\) 上的双射。

设 \(i \in I\) 与 \(x \in M_i\) ；由于 \(\Psi_i\) 是双射，一个映射 \(f_{i,x}\) （从 \(\Sigma\) 到自身）由下式定义：

\[
f _ {i, x} (\Psi_ {i} (\xi , \sigma)) = \Psi_ {i} (x \xi , \sigma) \quad (\xi \in \mathbf {M} _ {i}, \sigma \in \Sigma_ {i}).\tag{6}
\]

进一步，对于 \(a \in \mathbf{A}, f_a\) ，用它表示由下式定义的 \(\Sigma\) 到自身的映射：

\[
f _ {a} (a ^ {\prime}; i _ {1}, \dots , i _ {n}; p _ {1}, \dots , p _ {n}) = (a a ^ {\prime}; i _ {1}, \dots , i _ {n}; p _ {1}, \dots , p _ {n}).\tag{7}
\]

显然 \(f_{i, e_i}\) 是 \(\Sigma\) 的恒等映射，并且 \(f_{i, xx'} = f_{i, x} \circ f_{i, x'}\) 对 \(x, x'\) （在 \(M_i\) 中）成立，而 \(f_{i, h_i(a)} = f_a\) 对 \(a \in A\) 与 \(i \in I\) 成立。

于是命题 4 可应用于下述情形： \(\mathbf{M}'\) 是 \(\Sigma\) 到自身的映射所成的幺半群，其合成律为 \((f, f') \mapsto f \circ f'\) ，而 \(f_{i}\) 是同态 \(x \mapsto f_{i,x}\) ，它把 \(\mathbf{M}_{i}\) 映入 \(\mathbf{M}'\) ；则存在一个同态 \(f\) ，它把 \(\mathbf{M}\) 映入 \(\mathbf{M}'\) ，使得 \(f_{i,x} = f(\phi_i(x))\) 对 \(i \in \mathbf{I}\) 与 \(x \in \mathbf{M}_i\) 成立。设

\[
\sigma = (a; i _ {1}, \dots , i _ {n}; p _ {1}, \dots , p _ {n})
\]

为 \(\Sigma\) 中的序列。公式 (5) 至 (7) 通过对 \(n\) 作归纳给出关系

\[
\begin{array}{l} \sigma = (f _ {a} \circ f _ {i _ {1}, p _ {1}} \circ \dots \circ f _ {i _ {n}, p _ {n}}) (e) \\ = f (h (a) \phi_ {i _ {1}} (p _ {1}) \dots \phi_ {i _ {n}} (p _ {n})) (e), \end{array}
\]

即 \(\sigma = f(\Phi(\sigma))(e)\) 。这就证明了 \(\Phi\) 是单射。

\BourbakiRunInHeading{(B) 分解的存在性：}

设 D 为 M 中容许有分解的元素所成的集合。则 \(e \in D\) ，且 M 由 \(\bigcup_{i \in I} \phi_i(M_i)\) 、从而由 \(h(A) \cup \bigcup_{i \in I} \phi_i(P_i)\) 生成。于是 \(D \cdot \phi_i(P_i) \subset D\) 对所有 \(i \in I\) 成立；因此要证明 D = M，只需证明关系式 \(D \cdot h(A) \subset D\) ，而这由下面更精确的引理得出：

引理 1. 设 \(i_1, \ldots, i_n\) 属于 I，且 \(p_\alpha\) 在 \(\mathbf{P}_{i_\alpha}\) 中（\(1 \leqslant \alpha \leqslant n\) ）。对所有 \(a \in A\) ，存在 \(a' \in A\) 及一个序列 \((p'_\alpha)_{1 \leqslant \alpha \leqslant n}\) ，其中 \(p'_\alpha \in \mathbf{P}_{i_\alpha}\) ，使得

\[
\phi_ {i _ {1}} (p _ {1}) \dots \phi_ {i _ {n}} (p _ {n}) h (a) = h (a ^ {\prime}) \phi_ {i _ {1}} (p _ {1} ^ {\prime}) \dots \phi_ {i _ {n}} (p _ {n} ^ {\prime}).
\]

\(h(a) = \phi_{i_n}(h_{i_n}(a))\) ，并且存在 \(a_{n} \in \mathbf{A}\) 与 \(p_n' \in \mathbf{P}_{i_n}\) 满足

\[
p _ {n} \cdot h _ {i _ {n}} (a) = h _ {i _ {n}} (a _ {n}) \cdot p _ {n} ^ {\prime}.
\]

由此可得 \(\phi_{i_n}(p_n)h(a) = h(a_n)\phi_{i_n}(p_n')\) ，从而

\[
\phi_ {i _ {1}} (p _ {1}) \dots \phi_ {i _ {n - 1}} (p _ {n - 1}) \phi_ {i _ {n}} (p _ {n}) h (a) = \phi_ {i _ {1}} (p _ {1}) \dots \phi_ {i _ {n - 1}} (p _ {n - 1}) h (a _ {n}) \phi_ {i _ {n}} (p _ {n} ^ {\prime});
\]

引理由此通过对n的归纳得出。

\BourbakiRunInHeading{(C) 证明结束：}

设 \(x \in M\) ，n 为 x 的诸分解长度的最小值。我们将证明 x 的每个长度为 n 的分解 \(\sigma\) 都是既约的。这就确立了对 x 的既约分解的存在性；既约分解的唯一性进而蕴含 \(l(\sigma') > l(\sigma)\) 对 x 的每个分解 \(\sigma' \neq \sigma\) 成立。

情形 \(n = 0\) 是平凡的，故设 \(n > 0\) 。设

\[
\sigma = (a; i _ {1}, \dots , i _ {n}; p _ {1}, \dots , p _ {n})
\]

为 \(x\) 的一个长度为 \(n\) 的分解。如果存在整数 \(\alpha\) 满足 \(1 \leqslant \alpha \leqslant n\) 且 \(p_{\alpha} = e_{i_{\alpha}}\) ，则序列

\[
(a; i _ {1}, \dots , i _ {\alpha - 1}, i _ {\alpha + 1}, \dots , i _ {n}; p _ {1}, \dots , p _ {\alpha - 1}, p _ {\alpha + 1}, \dots , p _ {n})
\]

将是 \(x\) 的一个长度为 \(n - 1\) 的分解，这是被排除的。假设存在整数 \(\alpha\) 满足 \(1 \leqslant \alpha < n\) 且 \(i_{\alpha} = i_{\alpha + 1}\) ，并设

\[
p _ {\alpha} p _ {\alpha + 1} = h _ {i _ {\alpha}} \left(a ^ {\prime}\right) \cdot p _ {\alpha} ^ {\prime}
\]

其中 \(a' \in \mathbf{A}\) , \(p_{\alpha}' \in \mathrm{P}_{t_{\alpha}}\) ；根据引理 1 ，存在元素 \(a'' \in \mathbf{A}\) ,

\[
p _ {1} ^ {\prime} \in \mathrm{P} _ {i _ {1}}, \dots , p _ {\alpha - 1} ^ {\prime} \in \mathrm{P} _ {i _ {\alpha - 1}}
\]

使得

\[
\phi_ {i _ {1}} (p _ {1}) \dots \phi_ {i _ {\alpha - 1}} (p _ {\alpha - 1}) h (a ^ {\prime}) = h (a ^ {\prime \prime}) \phi_ {i _ {1}} (p _ {1} ^ {\prime}) \dots \phi_ {i _ {\alpha - 1}} (p _ {\alpha - 1} ^ {\prime})
\]

且序列

\[
\left(a a ^ {\prime \prime}; i _ {1}, \dots , i _ {\alpha - 1}, i _ {\alpha}, i _ {\alpha + 2}, \dots , i _ {n}; p _ {1} ^ {\prime}, \dots , p _ {\alpha - 1} ^ {\prime}, p _ {\alpha} ^ {\prime}, p _ {\alpha + 2}, \dots , p _ {n}\right)
\]

是 \(x\) 的一个长度为 \(n - 1\) 的分解，这是一个矛盾。

这样我们就证明了 \(\sigma\) 是既约的。

推论。在假设 (A) 下，同态 \(\phi_i\) 与 \(h\) 都是单射。对 I 中 \(i \neq j\) ，有 \(\phi_i(M_i) \cap \phi_j(M_j) = h(A)\) 。

首先 \(h\) 是单射的：如果 \(h(a) = h(a')\) ，则 (a) 与 \((a')\) 是 \(\mathbf{M}\) 中同一元素的两个既约分解，由此 \(a = a'\) 。设 \(i \in \mathbf{I}\) ；于是 \(h(\mathbf{A}) = \phi_i(h_i(\mathbf{A})) \subset \phi_i(\mathbf{M}_i)\) ；既约分解的唯一性蕴含

\[
h (\mathrm{A}) \cap \phi_ {i} (\mathrm{M} _ {i} - h _ {i} (\mathrm{A})) = \varnothing ,
\]

，从而 \(\phi_{i}(\mathbf{M}_{i} - h_{i}(\mathbf{A})) = \phi_{i}(\mathbf{M}_{i}) - h(\mathbf{A}).\)

诸同态 \(\phi_{i}\) 的单射性以及关系

\[
\phi_ {i} (\mathbf {M} _ {i}) \cap \phi_ {j} (\mathbf {M} _ {j}) \subset h (\mathrm{A})
\]

对 \(i \neq j\) 成立这一点则是下述事实的推论：对 \(i, j\) 属于 \(I, x\) 属于 \(M_i = h_i(A)\) 与 \(y\) 属于 \(M_j = h_j(A)\) ，关系 \(\phi_i(x) = \phi_j(y)\) 蕴含 \(i = j\) 且 \(x = y\) 。设 \(x = h_i(a).p\) 与 \(y = h_j(b).q\) ，其中 \(a, b\) 在 \(A, p\) 在 \(P_i = \{e_i\}\) 中且 \(y\) 在 \(P_j = \{e_j\}\) 中。则 \(\phi_i(x) = h(a)\phi_i(p)\) 且 \(\phi_j(y) = h(b)\phi_j(q)\) ，因此 \((a; i; p)\) 与 \((b; j; q)\) 是 \(M\) 中同一元素的两个既约分解。由此可得 \(i = j, a = b\) 且 \(p = q\) ，从而 \(x = h_i(a)p = h_j(b)q = y\) 。

当假设 (A) 成立时，我们将把每个幺半群 \(M_{i}\) 借助 \(\phi_{i}\) 与 M 的一个子幺半群等同起来；类似地，把 A 借助 h 与 M 的一个子幺半群等同起来。于是 M 由 \(\bigcup_{i\in I}M_{i}\) 生成，且 \(M_{i}\cap M_{j}=A\) 对 \(i\neq j\) 成立。

M 的每个元素都可以唯一地写成 \(a.\prod_{\alpha=1}p_{\alpha}\) 的形式，其中 \(a\in A, p_{1}\in P_{i1}=\{e\},\ldots,p_{n}\in P_{in}=\{e\}\) 且 \(i_{\alpha}\neq i_{\alpha+1}\) 对 \(1\leqslant\alpha<n\) 成立。最后，若 \(M'\) 是一个幺半群，且 \((f_{i}:M_{i}\to M')\) （\(i\in I\) ）是一族同态，它们在 A 上的限制都是 A 到 \(M'\) 中的同一个同态，则存在且仅存在一个同态 \(f:M\to M'\) ，它诱导出 \(f_{i}\) （在 \(M_{i}\) 上），对所有 \(i\in I\) 。

假设 (A) 在两种重要情形下成立：

\begin{enumerate}
\def\labelenumi{(\alph{enumi})}
\tightlist
\item
  \(\mathbf{A} = \{e\}\) 。在这种情形，有一族幺半群 \((\mathbf{M}_i)_{i\in I}\) ，而 \(\mathbf{M}\) 称为该族的幺半和。每个 \(\mathbf{M}_i\) 与 \(\mathbf{M}\) 的一个子幺半群等同，且 \(\mathbf{M}\) 由 \(\bigcup_{i\in I}\mathbf{M}_i\) 生成；进一步，\(\mathbf{M}_i\cap \mathbf{M}_j = \{e\}\) 对 \(i\neq j\) 成立。 \(\mathbf{M}\) 的每个元素都可以唯一地写成 \(x_{1}\ldots x_{n}\) 的形式，其中
\end{enumerate}

\[
x _ {1} \in \mathrm{M} _ {i _ {1}} - \{e \}, \dots , x _ {n} \in \mathrm{M} _ {i _ {n}} - \{e \}
\]

且 \(i_{\alpha} \neq i_{\alpha + 1}\) （\(1 \leqslant \alpha \leqslant n\) ）。最后，对每一族同态 \((f_i: \mathbf{M}_i \to \mathbf{M}')\) ，存在唯一的同态 \(f: \mathbf{M} \to \mathbf{M}'\) ，它在 \(\mathbf{M}_i\) 上的限制为 \(f_i\) ，对所有 \(i \in I\) 。

\begin{enumerate}
\def\labelenumi{(\alph{enumi})}
\setcounter{enumi}{1}
\tightlist
\item
  存在一族群 \((\mathbf{G}_i)_{i\in I}\) ，它们包含同一个群 A 作为子群，且 \(h_i\) 是 A 到 \(\mathbf{G}_i\) 中的单射。族 \((\mathbf{G}_i)_{i\in I}\) 由 A 融合后的和这时是一个群 G：幺半群 G 由 \(\bigcup_{i\in I}\phi_i(\mathbf{G}_i)\) 生成，且 \(\bigcup_{i\in I}\phi_i(\mathbf{G}_i)\) 的每个元素在 G 中都有逆元（参见 § 2，第 3 号，命题 4 的推论 1）；它记作 \(\star_{\mathrm{A}}\mathrm{G}_{i}\) ，也记作 \(\mathrm{G}_{1}*\mathrm{A}\mathrm{G}_{2}\) （当 \(\mathbf{I} = \{1,2\}\) ）。当 A 仅由单位元组成时，也称 G 是群族 \((\mathbf{G}_i)_{i\in I}\) 的自由积，记作 \(\star \mathbf{G}_i\) （或 \(\mathbf{G}_1*\mathbf{G}_2\) ，若 \(\mathbf{I} = \{1,2\}\) ）。†
\end{enumerate}

